% ============================================================
%  Section 1: Introduction
% ============================================================

Nonlinear model predictive control (NMPC) is the natural candidate for exothermic reactors: it predicts with the Arrhenius kinetics instead of a local linearization, it handles valve amplitude and rate constraints explicitly, and modern tools such as CasADi/IPOPT and do-mpc solve the underlying optimal control problems in milliseconds \cite{rawlings2020model, hedengren2014nonlinear, fiedler2023dompc}. Fixed-gain PID control, by contrast, still regulates the large majority of industrial loops \cite{astrom2006pid, seborg2010process}. The engineering question is therefore not whether NMPC can beat PID, but \emph{when} the additional model, estimator and on-line optimization pay off over a PID loop that has been tuned as well as a practitioner can tune it.

Published comparisons rarely answer this question, for two reasons. First, the linear baseline is often a classical empirical tuning correlation derived for unconstrained linear plants and known to give small robustness margins \cite{odwyer2009handbook, garpinger2014performance}. Beating such a baseline shows that robustness matters, not that nonlinear prediction does. Robust rules that bound the maximum sensitivity $M_s$, such as AMIGO \cite{astrom2006pid}, SIMC \cite{skogestad2003simple, grimholt2018optimal} or direct optimization of a time-domain criterion under an $M_s$ constraint \cite{garpinger2014performance, padula2012optimal}, are the relevant references. Second, the PID loop is frequently simulated without the actuator the NMPC is constrained by, or with an anti-windup scheme that ignores the valve slew rate, so that integrator windup rather than the control law decides the comparison \cite{hippe2006windup, tarbouriech2009antiwindup, bohn1995analysis, peng1996nonlinear}. A fair benchmark gives both controllers the same saturating, rate-limited valve, the same measurements and the same evaluation protocol.

Exothermic continuous stirred tank reactors (CSTRs) are the canonical test bed for this comparison \cite{uppal1974dynamic, aris1958analysis, luyben2007chemical, bequette2003process}: the exponential temperature dependence of the reaction rate makes the process gain vary with the operating point, and jacket fouling reduces the heat-transfer capacity during operation. Nonlinear geometric and robust designs have been proposed for this class of process \cite{alvarez1995robust}, but their benefit over a robustly tuned PID acting through a rate-limited valve has not been quantified scenario by scenario.

This paper performs that quantification on a three-state CSTR with an explicit jacket energy balance and a cooling valve limited to $[0, 300]$~L/min and $80$~L/min$^2$. The NMPC is evaluated in two forms: with access to the true state (NMPC-Oracle), which bounds what nonlinear prediction can achieve, and with an augmented moving-horizon estimator that reconstructs the concentration and a fouling factor from two thermocouples (NMPC-MHE), which is the implementable controller. The linear baselines are the two best robust tunings of a broader benchmark of the same plant: AMIGO and a PID whose ITAE is minimized subject to $M_s \le 1.6$ on the identified dead-time model. Both PID loops use a directional anti-windup that freezes the integrator whenever the previous command was cut by the amplitude \emph{or} the slew-rate limit.

The contributions are:
\begin{enumerate}
    \item A controlled comparison of NMPC and robust PID on the same nonlinear plant, valve and measurements, over six scenarios: nominal servo step, unmeasured feed disturbances, abrupt and graded jacket fouling, wide-range multi-point tracking and Monte Carlo sensor noise.
    \item A formulation of output-feedback NMPC with an augmented state $[C_A, T, T_j, \alpha_{UA}]^\top$, estimated by an extended Kalman filter whose estimates serve as arrival-cost priors for a moving-horizon refinement, together with a study of its sensitivity to the estimator covariances and of its computational cost per sample.
    \item Five criteria, each supported by a specific scenario, that delimit when nonlinear control pays off on this reactor and when a robust PID is the better choice, including a structural explanation of the offset of NMPC under disturbances absent from its model.
\end{enumerate}
All tables and figures are generated by the released simulation scripts with fixed random seeds; no number in this paper is entered by hand.

The paper is organized as follows. Section~\ref{sec:formulation} presents the plant, the NMPC and the augmented estimator, Section~\ref{sec:baselines} the PID baselines, Section~\ref{sec:results} the results and Section~\ref{sec:conclusion} the conclusions.
